% Source: https://artofproblemsolving.com/community/c3h1606259_binomial_coefficients
% Stable identity: aops-topic-1606259-q01
\textbf{(a)} Count the ordered lists of $k$ different elements of the $n$-element set in two ways.

Directly: the first entry can be chosen in $n$ ways, the second in $n-1$ ways, and so on, the $k$-th in $n-k+1$ ways, so there are
\[
n(n-1)\cdots(n-k+1)=\dfrac{n!}{(n-k)!}
\]
lists. Through subsets: every list is obtained by first choosing the set of its $k$ entries (in $\dbinom nk$ ways) and then putting these $k$ elements in order (in $k!$ ways), and different choices give different lists. So the number of lists is also $\dbinom nk\cdot k!$. Equating,
\[
\dbinom nk\cdot k!=\dfrac{n!}{(n-k)!}\qquad\Longrightarrow\qquad\dbinom nk=\dfrac{n!}{k!\,(n-k)!}.
\]
(For $k=0$ both sides are $1$: the empty set is the only $0$-element subset, and $0!=1$.)

\textbf{(b)} Write $(x+y)^n=(x+y)(x+y)\cdots(x+y)$ with $n$ factors, numbered $1$ to $n$. Multiplying out (by the distributive law) gives a sum of $2^n$ products, one for each way of choosing either $x$ or $y$ from every factor. A choice produces $x^ky^{n-k}$ exactly when $x$ is taken from $k$ of the factors and $y$ from the others, and such a choice is determined by the set of $k$ factors that contribute $x$, a $k$-element subset of $\{1,\dots,n\}$. Hence $x^ky^{n-k}$ appears exactly $\dbinom nk$ times, and
\[
(x+y)^n=\sum_{k=0}^{n}\dbinom nk x^ky^{n-k}.\qquad\blacksquare
\]
For example, in $(x+y)^3=x^3+3x^2y+3xy^2+y^3$ the term $x^2y$ comes from the three choices $xxy$, $xyx$, $yxx$, matching the $\dbinom32=3$ two-element subsets $\{1,2\},\{1,3\},\{2,3\}$ of factors.
